\documentclass[UTF8,a4paper,11pt]{ctexart}

\usepackage{amsmath,amssymb,bm}
\usepackage{geometry}
\usepackage{graphicx}
\usepackage{booktabs}
\usepackage{array}
\usepackage{float}
\usepackage{caption}
\usepackage{subcaption}
\usepackage{enumitem}
\usepackage{xcolor}
\usepackage{hyperref}

\geometry{left=2.4cm,right=2.4cm,top=2.5cm,bottom=2.5cm}

% ---------------------------------------------------------------------------
% 插图占位框：后续可用 \includegraphics 直接替换 \figbox 即可。
% 用法：\figbox{高度}  或  \figbox[宽度]{高度}
% ---------------------------------------------------------------------------
\newcommand{\figbox}[2][0.82\textwidth]{%
  \setlength{\fboxsep}{0pt}%
  \fcolorbox{gray!55}{gray!7}{%
    \begin{minipage}[c][#2][c]{#1}%
      \centering
      {\itshape\color{gray!65} 【插图占位框 \quad 后续替换为实际图片】}
    \end{minipage}}}

\title{\heiti 菲涅耳波带片的目标光场编码设计：\\稳妥基准方案与目标光场编码创新方案}
\author{项目组名称}
\date{\today}

\begin{document}

\maketitle

\begin{abstract}
菲涅耳波带片通常被作为一种固定焦距的衍射透镜来研究，其基本实验往往是根据给定焦距设计环带半径，然后验证其聚焦能力。然而，对于大学生物理实验竞赛而言，仅仅证明“波带片能够聚焦”创新性不足。本项目尝试从波动光学和光场调控的角度重新理解菲涅耳波带片：它不仅是一个衍射透镜，也可以被看作由透明/不透明区域组成的二值光场编码器。

为兼顾可实现性与创新性，本项目将研究内容组织为两个相互衔接的方案。\textbf{方案一（制备约束下的稳妥基准方案）}完整实现标准圆形菲涅耳波带片的设计、仿真、制备、聚焦测量、成像验证、色差测试与不确定度分析，并比较圆形、线型与方形等不同形状波带片的衍射聚焦特性，确保项目满足竞赛基本要求并提供可靠对照基准。\textbf{方案二（目标光场编码创新方案）}在方案一的平台与对照组之上展开：先给定目标轴向光强分布，再设计相应的二值波带结构，使实际衍射光场尽可能接近目标光场；其内部又分为\textbf{方案二-A：确定性扇区编码}与\textbf{方案二-B：径向二值优化}两条具体路线，分别实现单焦点、双焦点与扩展焦深光场。

两方案并非并列的两个研究方向，而是“基准—创新”的递进关系。核心创新在于设计思想由传统的“固定焦距正向设计”转变为“目标光场驱动的编码设计”。预期结果将揭示峰值强度、焦深、旁瓣与成像鲁棒性之间的物理权衡关系。
\end{abstract}

\noindent\textbf{关键词：} 菲涅耳波带片；衍射聚焦；二值编码；目标光场；双焦点；扩展焦深；物理实验创新

% ===========================================================================
\section{项目定位与核心创新}

\subsection{传统实验的不足}

菲涅耳波带片是一类典型的衍射光学元件。传统设计通常从一个给定焦距 $f$ 出发，根据相邻半波带到焦点的光程差条件计算环带半径。实验任务往往包括图案制作、光路搭建、焦距测量和误差分析。这类实验能够较好地验证惠更斯--菲涅耳原理和光的衍射干涉，但从创新性角度看，容易停留在“按公式制作一个波带片并观察聚焦”的层面。

本项目希望避免这种单纯验证式实验，而是将问题重新表述为：

\begin{quote}
\textbf{能否把菲涅耳波带片看作一个二值光场编码器，并通过编码方式主动设计其轴向光强分布？}
\end{quote}

在这种表述下，波带片不再只是一个固定焦距透镜，而是一个可以控制光场分布的二值衍射结构。实验的重点也从“焦距是否等于理论值”转向“实际光场是否接近设计目标”。

\subsection{两级方案的总体结构}

为兼顾“可实现性”与“创新性”，本项目将研究内容组织为两个相互衔接、层层递进的方案：

\begin{itemize}[leftmargin=2.4em]
\item \textbf{方案一（制备约束下的稳妥基准方案）：} 完整实现标准菲涅耳波带片的设计、仿真、制备、聚焦、成像、色差与不确定度分析，并比较不同形状波带片的聚焦特性。它确保项目满足竞赛基本要求，并为方案二提供实验平台与单焦点对照组。
\item \textbf{方案二（目标光场编码创新方案）：} 将波带片重新理解为二值光场编码器，由目标轴向光强反向设计二值结构。其内部又分为两条具体路线：\textbf{方案二-A：确定性扇区编码}与\textbf{方案二-B：径向二值优化}。
\end{itemize}

需要强调的是，方案二-A 与方案二-B 不是整个项目的两个并列大方向，而是“目标光场编码创新方案”内部的两个子方法；方案一与方案二之间则是“基准—创新”的递进关系：方案一保证稳妥与完整，方案二在其基础上展开光场编码创新。本报告后续先介绍两方案共用的基础物理模型，再分别展开方案一与方案二，最后给出统一的实验、不确定度与时间安排。

% ===========================================================================
\section{基础物理模型}

\subsection{标准菲涅耳波带片半径公式}

设平面波沿 $z$ 轴入射，目标焦点位于光轴上 $z=f$ 处。孔径面上半径为 $\rho$ 的点到焦点的距离为

\[
L(\rho)=\sqrt{\rho^2+f^2}.
\]

中心点到焦点的距离为 $f$，因此光程差为

\[
\Delta(\rho)=\sqrt{\rho^2+f^2}-f.
\]

相邻半波带的边界满足

\[
\Delta(r_n)=\frac{n\lambda}{2},
\]

其中 $\lambda$ 为入射光波长，$n=1,2,3,\ldots$ 为波带边界序号。由此得到

\[
\sqrt{r_n^2+f^2}-f=\frac{n\lambda}{2},
\]

整理可得精确半径公式：

\[
r_n=\sqrt{n\lambda f+\frac{n^2\lambda^2}{4}}.
\]

当 $f\gg n\lambda$ 时，常用近轴近似为

\[
r_n\approx \sqrt{n\lambda f}.
\]

标准二值振幅型波带片通常令相邻半波带交替透光和遮光，使到焦点处的透光波带尽可能相干叠加，其结构见图~\ref{fig:fzp-structure}。

\begin{figure}[H]
\centering
\includegraphics[width=0.52\textwidth]{figures/fig2_fzp_structure.pdf}
\caption{标准菲涅耳波带片的环带结构与半径定义：交替的透光/遮光半波带及其边界半径 $r_n$。}
\label{fig:fzp-structure}
\end{figure}

\subsection{二值透过率函数}

定义针对焦距 $f$ 的半波带序号函数

\[
q_f(\rho)=
\left\lfloor
\frac{2\left(\sqrt{\rho^2+f^2}-f\right)}{\lambda}
\right\rfloor ,
\]

其中 $\lfloor \cdot \rfloor$ 表示向下取整。二值振幅波带片的透过率可写成

\[
T_f(\rho)=
\begin{cases}
1, & q_f(\rho)\ \text{为偶数},\\
0, & q_f(\rho)\ \text{为奇数}.
\end{cases}
\]

这里 $T_f=1$ 表示透光，$T_f=0$ 表示遮光。若将偶数和奇数对调，只会改变选通的半波带组，不影响本项目的核心分析。

\subsection{菲涅耳传播模型}

设波带片位于 $z=0$ 平面，入射场为均匀平面波。波带片后的复振幅为

\[
U(\rho,\theta,0)=U_0 T(\rho,\theta),
\]

其中 $T(\rho,\theta)$ 是待设计的二值透过率函数。采用菲涅耳近似，传播到距离 $z$ 处轴上点的复振幅可写为

\[
U(0,0,z)=
\frac{e^{ikz}}{i\lambda z}
\int_0^{2\pi}
\int_0^R
T(\rho,\theta)
\exp\left(i\frac{k\rho^2}{2z}\right)
\rho\,d\rho\,d\theta,
\]

其中

\[
k=\frac{2\pi}{\lambda}
\]

为波数，$R$ 为波带片外半径。轴向光强为

\[
I(z)=|U(0,0,z)|^2.
\]

在实验和仿真中，常将轴向光强归一化为

\[
\widehat I(z)=\frac{I(z)}{\max_z I(z)}.
\]

本项目关注的不仅是单个焦点位置，更是整个轴向光强曲线 $\widehat I(z)$ 是否符合预设目标。

\subsection{制备约束与基准参数}

菲涅耳波带片的性能不仅由理论公式决定，还受到最小可制备线宽的限制。对标准圆形波带片，最外圈线宽近似为

\[
\Delta r_N \approx \frac{\lambda f}{2R}.
\]

因此，在给定波长 $\lambda$、外半径 $R$ 和最小线宽 $\Delta r_{\min}$ 时，焦距不能任意选择，而应满足

\[
f \approx \frac{2R\Delta r_N}{\lambda}.
\]

波带数与数值孔径近似为

\[
N\approx \frac{R^2}{\lambda f},\qquad
NA\approx \frac{R}{f}.
\]

这意味着更大的孔径有利于减小焦斑，但会导致更小的最外圈线宽，从而提高制备难度。本项目两方案统一采用如下参数作为设计基准：

\[
\lambda=532\,\mathrm{nm},\qquad
R=2\,\mathrm{mm},\qquad
f_0=150\,\mathrm{mm}.
\]

此时最外圈线宽约为 $\Delta r_N\approx 20\,\mu\mathrm{m}$，波带数约为 $N\approx 50$。这一参数组的优点是焦距适合桌面光路，最小线宽在高分辨率菲林、光绘掩膜或基础微加工条件下具有可实现性。

% ===========================================================================
\section{方案一：制备约束下的稳妥基准方案}

\subsection{方案定位}

在开展目标光场编码之前，本项目首先需要完成一个稳妥可靠的基准方案。该方案的目标不是追求复杂的光场调控，而是完整实现菲涅耳波带片的设计、仿真、制备、聚焦测量、成像验证和不确定度分析。它对应比赛题目中的基本要求，也是后续创新编码方案的实验基础和对照组。

基准方案的核心问题是：

\begin{quote}
\textbf{在实际制备精度有限的条件下，如何设计并制作一个能够稳定聚焦和成像的菲涅耳波带片，并定量研究其聚焦性能及影响因素？}
\end{quote}

这一方案强调可实现性和完整性。通过该方案，即使后续创新编码波带片的实验效果不完全理想，项目仍然能够保证有清晰的实验主线和可靠的数据结果。

\subsection{标准圆形菲涅耳波带片设计}

标准菲涅耳波带片以单一焦点为设计目标。设入射光波长为 $\lambda$，目标焦距为 $f$，第 $n$ 个波带边界半径为 $r_n$。由相邻半波带到焦点的光程差条件可得

\[
r_n=\sqrt{n\lambda f+\frac{n^2\lambda^2}{4}}.
\]

在近轴条件下，若 $f\gg n\lambda$，可近似为

\[
r_n\approx \sqrt{n\lambda f}.
\]

本项目拟采用 $532\,\mathrm{nm}$ 绿光激光作为主要光源，并初步选取

\[
\lambda=532\,\mathrm{nm},\qquad
f=150\,\mathrm{mm},\qquad
R=2\,\mathrm{mm}.
\]

其中 $R$ 为波带片外半径。对应最外圈线宽近似为

\[
\Delta r_N\approx \frac{\lambda f}{2R}.
\]

代入上述参数可得

\[
\Delta r_N\approx 20\,\mu\mathrm{m}.
\]

波带数近似为

\[
N\approx \frac{R^2}{\lambda f}.
\]

该设计对应的波带数约为 $50$。这一参数选择兼顾了桌面光路可实现性和图案制备可行性：焦距不至于过长，最小线宽又处于高分辨率菲林、光绘掩膜或基础光刻工艺可能实现的范围内。

\subsection{制备约束驱动的参数选择}

普通设计往往先选定焦距和孔径，再根据公式计算波带半径。但在真实实验中，最小可制备线宽 $\Delta r_{\min}$ 往往是最关键的限制因素。因此，本项目采用“制备约束驱动设计”的思路。

由

\[
\Delta r_N\approx \frac{\lambda f}{2R}
\]

可得

\[
f\approx \frac{2R\Delta r_N}{\lambda}.
\]

这说明焦距、孔径和最外圈线宽不能独立选择。若制备方法只能稳定实现较大的线宽，则焦距必须相应增大，或者孔径必须减小。反之，若希望使用较短焦距和较大孔径，则必须提高图案制备分辨率。

因此，基准方案的第一项研究内容是：

\begin{quote}
\textbf{根据可实现的最小线宽反推合理的孔径、焦距和波带数，而不是任意指定波带片参数。}
\end{quote}

这一点虽然不如目标编码方案显得“炫”，但具有明确的实验物理意义，也能体现项目对制备误差和实验可行性的重视。

\subsection{不同形状波带片的对比}

比赛题目要求研究波带片的形状和特性。因此，基准方案不仅制作标准圆形菲涅耳波带片，还可以在相同设计波长、近似相同焦距和相同外尺寸条件下，对比不同形状的波带片。

拟比较的结构包括：

\begin{enumerate}[label=(\arabic*)]
\item 圆形菲涅耳波带片：作为标准基准，预期形成近似圆对称焦斑。
\item 线型菲涅耳波带片：只在一个方向上实现聚焦，预期形成焦线，类似柱面透镜。
\item 方形或正交型波带片：具有更明显的方向性衍射特征，可能形成十字形或方形旁瓣结构。
\end{enumerate}

三类结构的比较重点不是简单判断哪一种“更好”，而是研究不同空间对称性如何影响衍射场分布。圆形结构具有旋转对称性，能量更容易集中到轴上焦点；线型结构只在一个方向压缩光场，因此形成线状焦区；方形或正交结构具有直角对称性，可能在焦平面产生方向性旁瓣和十字状能量分布。这一部分可以作为“形状与聚焦特性”的稳定研究内容。

\begin{figure}[H]
\centering
\includegraphics[width=0.98\textwidth]{figures/fig7_shapes.pdf}
\caption{三种形状波带片及其预期焦区对比：圆形（圆对称焦斑）、线型（焦线）、方形/正交型（方向性旁瓣与十字状能量分布）。}
\label{fig:shapes}
\end{figure}

\subsection{仿真内容}

基准方案将使用菲涅耳衍射或角谱法进行数值仿真。设波带片透过率为 $T(x,y)$，入射场为均匀平面波，则波带片后的复振幅为

\[
U(x,y,0)=U_0T(x,y).
\]

传播到距离 $z$ 后的光场可表示为

\[
U(x,y,z)=\mathcal{P}_z\{U(x,y,0)\},
\]

其中 $\mathcal{P}_z$ 表示菲涅耳传播算符或角谱传播算符。光强为

\[
I(x,y,z)=|U(x,y,z)|^2.
\]

仿真将输出以下结果：

\begin{enumerate}[label=(\arabic*)]
\item 焦平面光强分布 $I(x,y,f)$；
\item 轴向光强分布 $I_{\max}(z)$；
\item 光束传播截面 $I(x,z)$；
\item 焦斑半高宽 $\mathrm{FWHM}$；
\item 不同形状波带片的旁瓣结构；
\item 不同波长下的焦距偏移。
\end{enumerate}

\subsection{制备方式}

基准方案采用二值振幅型波带片。透明区域透光，黑色区域遮光。可选制备方式包括：

\begin{enumerate}[label=(\arabic*)]
\item 高分辨率菲林或光绘掩膜；
\item 在透明基底上制作金属遮光图案；
\item 光刻或激光直写；
\item 条件有限时，使用较大焦距、较大线宽版本进行原理验证。
\end{enumerate}

制备完成后，需要尽量进行显微检查，获得实际最外圈线宽、图案缩放误差、边缘粗糙度和中心偏移等信息。这些参数将用于解释实验焦距与理论焦距之间的偏差。

\subsection{聚焦实验}

聚焦实验光路为

\[
\text{激光器}
\rightarrow
\text{扩束/准直}
\rightarrow
\text{波带片}
\rightarrow
\text{可移动 CMOS 相机}
\rightarrow
\text{图像采集与处理},
\]

其示意见图~\ref{fig:setup}。相机沿光轴方向移动，记录不同 $z$ 位置处的光强图像。对每一个位置提取峰值光强

\[
I_{\max}(z)=\max_{x,y}I(x,y,z),
\]

并计算焦斑尺寸

\[
\mathrm{FWHM}(z).
\]

实验焦距可由两种方法确定：

\[
f_1=\arg\max_z I_{\max}(z),
\]

\[
f_2=\arg\min_z \mathrm{FWHM}(z).
\]

最终可取

\[
f_{\mathrm{exp}}=\frac{f_1+f_2}{2}.
\]

若两种方法得到的焦距差异较大，说明焦斑形态、旁瓣或噪声对测量产生了影响，需要进一步分析。

\begin{figure}[H]
\centering
\includegraphics[width=0.95\textwidth]{figures/fig6_setup.pdf}
\caption{聚焦测试光路示意：激光器、扩束/准直系统、待测波带片、沿光轴可移动的 CMOS 相机及图像采集系统。该光路同时用于方案一与方案二的轴向光强测量。}
\label{fig:setup}
\end{figure}

\subsection{成像实验}

除聚焦测试外，基准方案还将验证菲涅耳波带片的成像能力。将透明字母、网格或分辨率板作为物体，使波带片在相机端形成图像。近似采用薄透镜成像关系：

\[
\frac{1}{u}+\frac{1}{v}=\frac{1}{f},
\]

其中 $u$ 为物距，$v$ 为像距。若取

\[
u=2f,
\]

则理论上有

\[
v=2f,
\]

此时成像放大率为

\[
M=\frac{v}{u}=1.
\]

需要注意，菲涅耳波带片存在零级光、多级焦点和明显色差，因此成像质量通常不如普通折射透镜。这个现象不是实验失败，而是波带片作为衍射光学元件的固有限制，应该作为实验讨论的一部分。

\subsection{色差实验}

对同一个波带片，焦距与波长满足近似关系

\[
f(\lambda)\approx \frac{r_n^2}{n\lambda}.
\]

因此

\[
f\propto \frac{1}{\lambda}.
\]

若在 $532\,\mathrm{nm}$ 绿光下设计焦距为 $150\,\mathrm{mm}$，则在红光和蓝光下焦距会发生明显移动。使用红光、绿光和蓝光激光分别测量焦距，可以直观展示衍射透镜的色差。

该实验也可以为后续目标编码方案提供参考：扩展焦深设计虽然不能真正消除色差，但可能在一定程度上提高对焦点偏移的容忍度。

\begin{figure}[H]
\centering
\includegraphics[width=0.9\textwidth]{figures/fig8_chromatic.pdf}
\caption{衍射色差示意：同一波带片在红、绿、蓝三波长下焦距按 $f\propto 1/\lambda$ 发生移动。}
\label{fig:chromatic}
\end{figure}

\subsection{不确定度分析}
\label{sec:baseline-uncertainty}

基准方案需要建立完整的不确定度分析框架。主要误差来源包括：

\begin{enumerate}[label=(\arabic*)]
\item 位移台或刻度尺的读数误差；
\item 相机像素尺寸标定误差；
\item 光斑拟合误差；
\item 波带片与光轴不垂直引起的对准误差；
\item 激光准直不完全；
\item 图案最外圈线宽误差；
\item 图案整体缩放误差；
\item 遮光区域不完全不透明；
\item 背景光和相机噪声。
\end{enumerate}

实验焦距不确定度可写为

\[
u(f_{\mathrm{exp}})
=
\sqrt{
u_z^2+
u_{\mathrm{fit}}^2+
u_{\mathrm{align}}^2
}.
\]

理论焦距由图案半径决定。由于

\[
f\approx \frac{r_n^2}{n\lambda},
\]

若图案存在半径误差 $\delta r_n$，则

\[
\frac{\delta f}{f}
\approx
2\frac{\delta r_n}{r_n}
-
\frac{\delta\lambda}{\lambda}.
\]

若图案存在整体缩放因子 $s$，即

\[
r_n'=sr_n,
\]

则实际焦距近似变为

\[
f'=s^2f.
\]

因此，图案只有 $1\%$ 的缩放误差，也可能导致约 $2\%$ 的焦距偏差。该结论可以作为解释实验误差的重要依据，并为方案二中编码波带片的焦点偏移分析提供基础。

\subsection{基准方案的作用与局限}

基准方案的作用包括：

\begin{enumerate}[label=(\arabic*)]
\item 确保项目能够完成比赛题目的基本要求；
\item 提供标准单焦点波带片作为后续编码波带片的对照组；
\item 建立仿真、制备、测试和不确定度分析流程；
\item 明确制备精度、最外圈线宽和实际聚焦性能之间的关系。
\end{enumerate}

但该方案的创新性相对有限。它主要回答“波带片能否按照理论焦距聚焦”和“形状如何影响焦斑”这类问题。为了进一步体现创新命题的特点，本项目在基准方案基础上提出目标光场编码方案，将波带片从固定焦距衍射透镜拓展为二值光场编码器。

% ===========================================================================
\section{方案二：目标光场编码创新方案}

\subsection{将菲涅耳波带片视为二值光场编码器}

方案二的出发点是改变对菲涅耳波带片的理解方式。传统波带片设计属于正向设计：

\[
f \longrightarrow r_n \longrightarrow T(x,y) \longrightarrow I(x,y,z),
\]

其中 $f$ 是目标焦距，$r_n$ 是环带边界半径，$T(x,y)$ 是波带片透过率函数，$I(x,y,z)$ 是传播后的光强分布。

方案二采用的思路更接近目标驱动设计：

\[
I_{\mathrm{target}}(z) \longrightarrow T(x,y) \longrightarrow I_{\mathrm{exp}}(z),
\]

即首先规定希望实现的轴向光强分布 $I_{\mathrm{target}}(z)$，再设计二值波带结构 $T(x,y)$，最后通过实验测量实际轴向光强 $I_{\mathrm{exp}}(z)$，并评价二者之间的差异（见图~\ref{fig:paradigm}）。这种思路的创新性不在于使用更复杂的设备，而在于改变了对菲涅耳波带片的理解方式：从“衍射透镜”转变为“二值光场编码器”。

在此框架下，方案二给出两条具体的编码路线：\textbf{方案二-A：确定性扇区编码}（§\ref{sec:sector}）以清晰的物理图像实现多焦点与扩展焦深；\textbf{方案二-B：径向二值优化}（§\ref{sec:binary}）以优化算法逼近任意目标轴向光场。两者共享同一目标光场定义（§\ref{sec:targets}），并在峰值强度、焦深与旁瓣之间体现出统一的物理权衡（§\ref{sec:tradeoff}）。

\begin{figure}[H]
\centering
\includegraphics[width=0.95\textwidth]{figures/fig1_paradigm.pdf}
\caption{正向设计（焦距驱动）与目标驱动设计（目标光场驱动）两种范式的对比示意图。}
\label{fig:paradigm}
\end{figure}

\subsection{目标轴向光强的定义}
\label{sec:targets}

方案二不只考虑标准单焦点，还设计双焦点和扩展焦深两类目标光场（见图~\ref{fig:targets}）。

\subsubsection{单焦点目标}

单焦点目标可用高斯函数表示：

\[
I_{\mathrm{target}}^{(1)}(z)
=
\exp\left[
-\frac{(z-f_0)^2}{2\sigma_z^2}
\right],
\]

其中 $f_0$ 为目标焦距，$\sigma_z$ 描述目标焦区宽度。标准菲涅耳波带片对应的就是这一类目标。

\subsubsection{双焦点目标}

双焦点目标可写为两个高斯峰的叠加：

\[
I_{\mathrm{target}}^{(2)}(z)
=
A_1
\exp\left[
-\frac{(z-f_1)^2}{2\sigma_z^2}
\right]
+
A_2
\exp\left[
-\frac{(z-f_2)^2}{2\sigma_z^2}
\right],
\]

其中 $f_1$ 和 $f_2$ 是两个目标焦距，$A_1$ 和 $A_2$ 控制两个焦点的相对强度。例如可取

\[
f_1=135\,\mathrm{mm},\qquad
f_2=165\,\mathrm{mm}.
\]

若希望两个焦点强度近似相等，可取 $A_1=A_2$。

\subsubsection{扩展焦深目标}

扩展焦深目标不是追求某一个位置具有极高峰值，而是希望在一段轴向范围内保持较高光强。可用超高斯函数表示：

\[
I_{\mathrm{target}}^{(3)}(z)
=
\exp\left[
-\left(\frac{z-f_0}{\Delta f}\right)^{2p}
\right],
\]

其中 $\Delta f$ 表示希望扩展的半焦深范围，$p$ 为超高斯阶数。$p$ 越大，目标曲线越接近平台型分布。例如可取

\[
f_0=150\,\mathrm{mm},\qquad
\Delta f=15\,\mathrm{mm},\qquad
p=4.
\]

扩展焦深设计体现了一个重要的物理权衡：峰值强度可能下降，但系统对离焦误差的容忍度提高。

\begin{figure}[H]
\centering
\includegraphics[width=0.98\textwidth]{figures/fig3_targets.pdf}
\caption{三类目标轴向光强分布示意：单焦点（高斯）、双焦点（双高斯峰）与扩展焦深（超高斯平台）。}
\label{fig:targets}
\end{figure}

% ---------------------------------------------------------------------------
\subsection{方案二-A：确定性扇区编码}
\label{sec:sector}

\subsubsection{基本思想}

方案二-A 的核心思想是：将圆形孔径划分为多个角向扇区，每个扇区仍然按照普通菲涅耳波带片公式设计，但不同扇区可以对应不同的焦距。这样，整个波带片可以看作多个子孔径的组合。每个子孔径倾向于在其设计焦距附近形成光强增强，多个子孔径叠加后便可形成双焦点或扩展焦深光场。

这种方法的优点是物理图像清楚、图案生成简单、可解释性强。它不是改变实验设备，而是改变同一个孔径内部的编码方式。

\subsubsection{数学表达}

将角向范围 $0\leq \theta <2\pi$ 划分为 $K$ 个扇区。第 $s$ 个扇区为

\[
\Theta_s=
\left[
\frac{2\pi(s-1)}{K},
\frac{2\pi s}{K}
\right),
\qquad s=1,2,\ldots,K.
\]

定义扇区编号函数

\[
s(\theta)=1+\left\lfloor \frac{K\theta}{2\pi}\right\rfloor .
\]

为每个扇区分配一个设计焦距 $f_s$。则扇区编码波带片的透过率可写为

\[
T_A(\rho,\theta)
=
T_{f_{s(\theta)}}(\rho),
\qquad 0\leq \rho \leq R,
\]

其中 $T_f(\rho)$ 是标准二值波带片透过率函数。一个 $K=8$ 的双焦点扇区编码示例见图~\ref{fig:sector}。

\begin{figure}[H]
\centering
\includegraphics[width=0.56\textwidth]{figures/fig4_sector.pdf}
\caption{确定性扇区编码示意：圆形孔径被划分为 $K$ 个角向扇区，奇偶扇区分别对应焦距 $f_1$ 与 $f_2$。}
\label{fig:sector}
\end{figure}

\subsubsection{三种编码方式}

\paragraph{单焦点基准。}
若所有扇区使用同一个焦距，

\[
f_s=f_0,\qquad s=1,2,\ldots,K,
\]

则退化为标准单焦点菲涅耳波带片。这一结构作为实验对照组。

\paragraph{双焦点扇区编码。}
若希望得到两个焦点，可令奇数扇区和偶数扇区对应不同焦距：

\[
f_s=
\begin{cases}
f_1, & s\ \text{为奇数},\\
f_2, & s\ \text{为偶数}.
\end{cases}
\]

例如取 $f_1=135\,\mathrm{mm}$，$f_2=165\,\mathrm{mm}$。这样，约一半孔径对第一个焦点相干贡献，另一半孔径对第二个焦点相干贡献。实验中预期在轴向光强曲线中出现两个主要峰值。

若希望两个焦点强度比为 $I_1:I_2$，可粗略按孔径面积比例调节分配给两个焦距的扇区数。由于焦点处的相干振幅近似与参与该焦点的孔径面积成正比，峰值强度近似与面积比例的平方相关。因此可用

\[
p_1:p_2\approx \sqrt{I_1}:\sqrt{I_2}
\]

作为初步设计，其中 $p_1$ 和 $p_2$ 分别为分配给两个焦距的孔径面积比例。实际结果还需通过菲涅耳传播仿真修正。

\paragraph{扩展焦深扇区编码。}
若希望得到扩展焦深，可将不同扇区的焦距分布在一个区间内：

\[
f_s \in [f_0-\Delta f, f_0+\Delta f].
\]

一种简单确定性分配方式为

\[
f_s=f_0-\Delta f+
\frac{2\Delta f}{K-1}(s-1),
\qquad s=1,2,\ldots,K.
\]

为了减小角向方向性衍射条纹，也可以不按相邻扇区连续排列焦距，而是将不同焦距的扇区交错分布，使整个孔径在角向上更均匀。

\subsubsection{物理解释}

标准波带片的所有透光波带都服务于同一个焦距，因此焦点峰值强，但焦深较短。双焦点编码将孔径资源分配给两个不同焦距，因此两个焦点都能形成，但每个焦点可利用的有效孔径减少，峰值强度降低，旁瓣和背景增强。扩展焦深编码将孔径资源分散到一段焦距范围内，结果是峰值不如标准波带片高，但轴向可用范围变长。

\subsubsection{方案二-A 的预期结果}

方案二-A 预计得到如下实验结果：

\begin{enumerate}[label=(\arabic*)]
\item 标准单焦点波带片的轴向光强曲线具有单一尖峰。
\item 双焦点扇区编码波带片在两个预设位置附近出现两个主峰。
\item 扩展焦深扇区编码波带片的轴向光强曲线峰值较低，但峰宽明显增加。
\item 在成像实验中，标准波带片在最佳像面处清晰度最高，但离焦后快速变差；扩展焦深波带片在最佳像面处锐度可能略低，但在较大轴向范围内保持可辨图像。
\end{enumerate}

% ---------------------------------------------------------------------------
\subsection{方案二-B：径向二值优化}
\label{sec:binary}

\subsubsection{基本思想}

方案二-B 比方案二-A 更进一步。它不再预设每个区域必须服从标准波带片公式，而是将波带片离散成若干径向环带，每个环带只有两个状态：透光或不透光。这样，波带片设计问题就转化为一个二值优化问题。

设第 $m$ 个环带的二值变量为

\[
a_m\in\{0,1\},
\]

其中 $a_m=1$ 表示第 $m$ 个环带透光，$a_m=0$ 表示第 $m$ 个环带遮光。设计任务变为寻找一组二值向量

\[
\bm a=(a_1,a_2,\ldots,a_M)
\]

使实际轴向光强 $\widehat I_{\bm a}(z)$ 尽可能接近目标光强 $I_{\mathrm{target}}(z)$。该方案的总体流程见图~\ref{fig:method-b}。

\begin{figure}[H]
\centering
\includegraphics[width=0.9\textwidth]{figures/fig5_method_b.pdf}
\caption{径向二值优化总览：目标轴向光强、二值环带变量 $a_m\in\{0,1\}$、由其生成的圆形二值波带片掩膜，以及目标驱动的优化迭代回路。}
\label{fig:method-b}
\end{figure}

\subsubsection{径向离散化}

将波带片半径 $0\leq \rho\leq R$ 划分为 $M$ 个环带：

\[
0=r_0<r_1<r_2<\cdots<r_M=R.
\]

为满足制备约束，需保证每个环带宽度不小于最小可制备线宽：

\[
r_m-r_{m-1}\geq \Delta r_{\min}.
\]

第 $m$ 个环带的指示函数定义为

\[
\Pi_m(\rho)=
\begin{cases}
1, & r_{m-1}\leq \rho<r_m,\\
0, & \text{其他}.
\end{cases}
\]

则径向二值优化波带片的透过率可写为

\[
T_B(\rho)
=
\sum_{m=1}^{M} a_m \Pi_m(\rho).
\]

与标准波带片相比，这里的 $a_m$ 不再由奇偶半波带直接决定，而是由目标光场优化得到。

\subsubsection{轴向场的离散表示}

将轴向采样点记为

\[
z_1,z_2,\ldots,z_J.
\]

第 $m$ 个环带对轴上点 $z_j$ 的复振幅贡献可定义为

\[
C_{mj}
=
\frac{2\pi e^{ikz_j}}{i\lambda z_j}
\int_{r_{m-1}}^{r_m}
\exp\left(i\frac{k\rho^2}{2z_j}\right)
\rho\,d\rho.
\]

因此，对于给定二值向量 $\bm a$，轴上点 $z_j$ 的总复振幅为

\[
U_{\bm a}(z_j)
=
\sum_{m=1}^{M} a_m C_{mj}.
\]

对应光强为

\[
I_{\bm a}(z_j)=
\left|
\sum_{m=1}^{M} a_m C_{mj}
\right|^2.
\]

归一化后得到

\[
\widehat I_{\bm a}(z_j)
=
\frac{I_{\bm a}(z_j)}
{\max_j I_{\bm a}(z_j)}.
\]

这一表达式揭示了方案二-B 的物理本质：每个环带都是一个向轴上不同位置贡献复振幅的“相量源”，优化过程就是选择哪些相量参与叠加，从而塑造轴向光强分布。

\subsubsection{目标函数}

为使实际光强逼近目标光强，可定义目标函数

\[
J(\bm a)
=
\sum_{j=1}^{J}
w_j
\left[
\widehat I_{\bm a}(z_j)
-
I_{\mathrm{target}}(z_j)
\right]^2
+
\alpha S_{\mathrm{sidelobe}}
+
\beta C_{\mathrm{edge}}
+
\gamma C_{\mathrm{min}}.
\]

其中第一项为目标匹配误差，$w_j$ 是不同轴向位置的权重。

$S_{\mathrm{sidelobe}}$ 用于惩罚目标焦区之外的光强，例如可定义为

\[
S_{\mathrm{sidelobe}}
=
\frac{
\sum_{z_j\in \Omega_{\mathrm{out}}}
\widehat I_{\bm a}(z_j)
}{
\sum_{j=1}^{J}
\widehat I_{\bm a}(z_j)
},
\]

其中 $\Omega_{\mathrm{out}}$ 表示目标焦区以外的轴向范围。

$C_{\mathrm{edge}}$ 用于惩罚过于频繁的开闭变化：

\[
C_{\mathrm{edge}}
=
\sum_{m=2}^{M}
|a_m-a_{m-1}|.
\]

这一项可以减少环带过度碎片化，使图案更容易制备。

$C_{\mathrm{min}}$ 用于惩罚小于最小可制备宽度的连续结构。若将连续相同状态的环带段称为一个 run，其长度为 $L_\ell$，则可定义

\[
C_{\mathrm{min}}
=
\sum_{\ell}
\left[
\max(0,L_{\min}-L_\ell)
\right]^2.
\]

其中 $L_{\min}$ 是允许的最小连续环带数。该项可以避免出现过窄的透光或遮光结构。

因此，方案二-B 的优化问题可写为

\[
\min_{\bm a\in\{0,1\}^{M}} J(\bm a).
\]

\subsubsection{优化算法}

本项目不需要使用复杂机器学习算法。可采用简单的二值局部搜索或模拟退火算法。基本流程如下：

\begin{enumerate}[label=\textbf{Step \arabic*:}]
\item 给定 $\lambda$、$R$、$\Delta r_{\min}$、轴向采样范围和目标光强 $I_{\mathrm{target}}(z)$。
\item 将半径方向离散成 $M$ 个环带，并预计算每个环带对每个轴向点的复振幅贡献 $C_{mj}$。
\item 初始化二值向量 $\bm a$。初始值可以来自标准菲涅耳波带片，也可以随机生成。
\item 随机选择一个环带 $m$，尝试翻转其状态：
\[
a_m\rightarrow 1-a_m.
\]
\item 计算新旧目标函数差值：
\[
\Delta J=J_{\mathrm{new}}-J_{\mathrm{old}}.
\]
\item 若 $\Delta J<0$，接受该变化；若采用模拟退火，也可以以概率
\[
P=\exp\left(-\frac{\Delta J}{T}\right)
\]
接受较差解，以避免陷入局部最优。
\item 逐步降低温度 $T$，重复迭代直到目标函数收敛。
\item 输出最终二值向量 $\bm a$，生成可制备的波带片图案。
\end{enumerate}

\subsubsection{方案二-B 的物理意义}

方案二-B 的创新点不在于算法复杂，而在于建立了明确的“目标光场--二值结构”映射。传统波带片的环带开闭由单焦点相干条件决定，而径向二值优化允许部分环带为双焦点、扩展焦深或旁瓣抑制服务。

从物理上看，轴上光场是所有透光环带贡献的相干叠加：

\[
U(z)=\sum_m a_m C_m(z).
\]

改变 $a_m$ 就是在选择哪些相量参与干涉。不同选择会使相干增强出现在不同的轴向位置。优化过程实际上是在利用干涉原理重新分配光能。

\subsubsection{方案二-B 的预期结果}

方案二-B 预计能够得到比方案二-A 更灵活的轴向光场。例如：

\begin{enumerate}[label=(\arabic*)]
\item 对双焦点目标，优化结果应在两个目标焦距附近形成双峰，并可通过目标函数控制两峰强度比例。
\item 对扩展焦深目标，优化结果应形成比标准波带片更宽的轴向强度平台。
\item 与方案二-A 相比，方案二-B 可能具有更好的目标匹配度，但图案解释性较弱，且可能对制备误差更敏感。
\item 通过加入 $C_{\mathrm{edge}}$ 和 $C_{\mathrm{min}}$ 等正则化项，可以在目标匹配和可制备性之间取得平衡。
\end{enumerate}

% ---------------------------------------------------------------------------
\subsection{峰值强度、焦深与旁瓣的权衡}
\label{sec:tradeoff}

无论是方案二-A 还是方案二-B，目标光场编码都受同一物理规律支配：孔径所能提供的相干能量是有限的，将其重新分配到多个焦点或一段焦距范围，必然以单点峰值强度的下降为代价。由此可得方案二的核心权衡关系：

\[
\text{高峰值强度}
\quad \Longleftrightarrow \quad
\text{短焦深},
\]

\[
\text{扩展焦深 / 多焦点}
\quad \Longleftrightarrow \quad
\text{峰值强度下降、旁瓣增加}.
\]

这正是本项目希望通过实验揭示的核心物理关系。方案二-A 与方案二-B 是实现同一编码目标的两条路线，二者的比较见表~\ref{tab:method-compare}。

\begin{table}[H]
\centering
\caption{方案二-A（确定性扇区编码）与方案二-B（径向二值优化）的比较}
\label{tab:method-compare}
\begin{tabular}{p{3.2cm}p{5.1cm}p{5.1cm}}
\toprule
比较项 & 方案二-A：确定性扇区编码 & 方案二-B：径向二值优化 \\
\midrule
核心思想 & 将孔径划分为多个角向扇区，不同扇区对应不同焦距 & 将径向环带作为二值变量，通过优化逼近目标光强 \\
物理图像 & 多个子波带片的孔径复用与相干叠加 & 多个环带相量的选择性干涉叠加 \\
设计自由度 & 中等，主要由扇区数和焦距分布决定 & 高，由所有环带开闭状态共同决定 \\
可解释性 & 强，容易向评委说明 & 较强，但需要解释优化目标和相量叠加 \\
图案复杂度 & 较规则，制备相对容易 & 可能较不规则，需要加入可制备性约束 \\
适合目标 & 双焦点、扩展焦深的直观演示 & 任意目标轴向光场的优化逼近 \\
潜在问题 & 角向扇区会带来方向性旁瓣和焦斑畸变 & 容易陷入局部最优，对制备误差敏感 \\
创新侧重点 & 用孔径复用实现多焦点和扩展焦深 & 将波带片设计转化为二值优化问题 \\
\bottomrule
\end{tabular}
\end{table}

两条路线并不互相排斥。方案二-A 适合作为直观、可靠的创新演示；方案二-B 适合作为更深入的目标编码设计。项目可以先完成方案二-A，再将方案二-B 作为进一步深化，从而形成由确定性设计到优化设计的完整编码研究路线。

% ===========================================================================
\section{统一实验方案}

方案一与方案二共用同一套激光光源、准直系统、二值振幅掩膜和可移动 CMOS 相机。本节给出覆盖两方案全部样品的统一实验与测量方案。

\subsection{样品设计}

拟设计并制作以下样品，其中标准与不同形状波带片属于方案一，编码波带片属于方案二：

\begin{table}[H]
\centering
\caption{拟制作的波带片样品}
\label{tab:samples}
\begin{tabular}{llll}
\toprule
样品编号 & 类型 & 设计目标 & 所属方案 \\
\midrule
S-FZP & 标准单焦点圆形波带片 & $f_0=150\,\mathrm{mm}$ & 方案一（基准/对照） \\
L-FZP & 线型波带片 & 焦线，$f_0\approx150\,\mathrm{mm}$ & 方案一（形状对比） \\
Q-FZP & 方形/正交型波带片 & 方向性旁瓣 & 方案一（形状对比） \\
D-Sector & 双焦点扇区编码波带片 & $f_1=135$, $f_2=165\,\mathrm{mm}$ & 方案二-A \\
E-Sector & 扩展焦深扇区编码波带片 & $135$--$165\,\mathrm{mm}$ & 方案二-A \\
D-Binary & 双焦点径向优化波带片 & 双峰目标光强 & 方案二-B \\
E-Binary & 扩展焦深径向优化波带片 & 平台型目标光强 & 方案二-B \\
\bottomrule
\end{tabular}
\end{table}

实际制作数量可根据制备条件调整。若时间有限，至少应完成方案一的标准与形状样品；若仿真和制备顺利，再加入方案二的编码样品。

\subsection{光路与制备}

聚焦与成像测试的光路、制备方式及制备质量检查均与方案一一致（见 §\ref{sec:baseline-uncertainty} 前的相应小节及图~\ref{fig:setup}）。CMOS 相机沿光轴在 $z=120$--$180\,\mathrm{mm}$ 范围内采集光强图像，粗扫步长可取 $1$--$2\,\mathrm{mm}$，在焦点附近细扫步长可取 $0.2$--$0.5\,\mathrm{mm}$。建议以 $532\,\mathrm{nm}$ 绿光激光为主，并在色差测试中加入红光与蓝光激光。

\subsection{数据处理}

对每个轴向位置 $z_j$，采集二维光强图像 $I(x,y;z_j)$。处理步骤包括：

\begin{enumerate}[label=(\arabic*)]
\item 扣除背景光；
\item 避免相机饱和；
\item 保持曝光时间和增益不变；
\item 提取峰值光强 $I_{\max}(z_j)$；
\item 提取焦斑半高宽 $\mathrm{FWHM}(z_j)$；
\item 计算旁瓣强度；
\item 与目标轴向光强进行比较。
\end{enumerate}

轴向峰值曲线定义为

\[
I_{\max}(z_j)=\max_{x,y} I(x,y;z_j),
\]

归一化后得到

\[
\widehat I_{\max}(z_j)=
\frac{I_{\max}(z_j)}
{\max_j I_{\max}(z_j)}.
\]

对方案二的编码波带片，进一步定义目标匹配误差

\[
E_{\mathrm{exp}}
=
\sum_j
\left[
\widehat I_{\max}(z_j)
-
I_{\mathrm{target}}(z_j)
\right]^2.
\]

该指标体现了方案二与普通焦距测量实验的不同：我们不是只判断焦点在哪里，而是评价整条轴向光强曲线是否符合设计目标。

\subsection{焦斑尺寸}

焦平面附近的光斑可用二维图像分析获得。对于中心截线，可用高斯函数近似拟合：

\[
I(x)=I_0
\exp\left[
-\frac{2(x-x_0)^2}{w^2}
\right]
+B.
\]

对应半高宽为

\[
\mathrm{FWHM}=w\sqrt{2\ln 2}.
\]

虽然实际菲涅耳波带片焦斑不一定严格为高斯型，但统一采用相同拟合方法可以保证不同样品之间的可比性。

\subsection{焦深定义}

为比较标准波带片和扩展焦深波带片，可定义有效焦深为满足以下条件的轴向范围长度：

\[
I_{\max}(z)\geq \eta I_{\max}^{\mathrm{peak}},
\]

其中 $\eta$ 可取 $0.5$ 或 $0.7$。也可同时要求焦斑尺寸不超过某一阈值：

\[
\mathrm{FWHM}(z)\leq \mathrm{FWHM}_{\mathrm{th}}.
\]

因此有效焦深可定义为

\[
DOF_{\eta}
=
\mathrm{length}
\left\{
z:
I_{\max}(z)\geq \eta I_{\max}^{\mathrm{peak}}
\right\}.
\]

\subsection{成像离焦容忍度}

为体现扩展焦深设计的应用意义，可进行简单成像实验。用透明字母、网格或分辨率板作为物体，使波带片形成图像。然后将相机放置在最佳像面及其前后不同位置，比较图像质量。图像对比度可定义为

\[
C=
\frac{I_{\max}-I_{\min}}
{I_{\max}+I_{\min}},
\]

也可用平均边缘梯度表示图像锐度：

\[
G=\left\langle |\nabla I(x,y)| \right\rangle.
\]

预期标准波带片在最佳像面处 $C$ 和 $G$ 较高，但离焦后迅速下降；扩展焦深波带片在最佳像面处可能略低，但随离焦距离变化更缓慢。

% ===========================================================================
\section{编码方案的目标匹配不确定度}

与焦距测量、制备误差有关的不确定度已在方案一的不确定度框架（§\ref{sec:baseline-uncertainty}）中系统讨论，此处不再重复。对方案二而言，由于关心的是整条轴向光强曲线，还应额外考虑目标匹配误差的不确定度。

若每个轴向位置重复测量 $M$ 次，则可计算

\[
\bar I(z_j)=
\frac{1}{M}
\sum_{m=1}^{M}
I_m(z_j),
\]

以及标准不确定度

\[
u_I(z_j)=
\sqrt{
\frac{1}{M(M-1)}
\sum_{m=1}^{M}
\left[
I_m(z_j)-\bar I(z_j)
\right]^2
}.
\]

目标匹配误差

\[
E_{\mathrm{exp}}
=
\sum_j
\left[
\widehat I_{\max}(z_j)
-
I_{\mathrm{target}}(z_j)
\right]^2
\]

的不确定度可通过重复实验或误差传播估算。这样可以使方案二的实验评价更加量化，而不是仅凭图像观察判断编码效果。

% ===========================================================================
\section{预期创新成果}

本项目预期形成以下成果：

\begin{enumerate}[label=(\arabic*)]
\item 在方案一中完成标准单焦点菲涅耳波带片的设计、仿真、制备与焦距测量，并比较圆形、线型与方形波带片的聚焦特性，作为实验基准与对照。
\item 在方案二-A 中通过确定性扇区编码实现双焦点和扩展焦深光场，展示孔径复用带来的多焦点调控能力。
\item 在方案二-B 中通过径向二值优化建立目标轴向光场与二值环带结构之间的联系，验证波带片作为二值光场编码器的可行性。
\item 通过实验比较标准单焦点、双焦点和扩展焦深波带片的轴向光强、焦斑尺寸、旁瓣和成像离焦容忍度。
\item 建立包含位移误差、拟合误差、像素标定误差、制备误差和目标匹配误差的不确定度分析框架。
\item 总结峰值强度、焦深、旁瓣和图像鲁棒性之间的物理权衡关系。
\end{enumerate}

本项目的创新性可概括为：

\begin{quote}
\textbf{不是研究“波带片能不能聚焦”，而是在稳妥基准之上研究“能否通过二值编码设计波带片的焦场”。}
\end{quote}

% ===========================================================================
\section{时间安排}

\begin{table}[H]
\centering
\caption{项目时间安排}
\label{tab:schedule}
\begin{tabular}{p{2cm}p{11cm}}
\toprule
阶段 & 工作内容 \\
\midrule
第 1 周 & 阅读菲涅耳波带片、衍射传播和二值光场编码相关资料，确定实验参数和目标光场。 \\
第 2 周 & 编写标准/形状波带片（方案一）与扇区编码、径向二值优化波带片（方案二）的图案生成程序。 \\
第 3 周 & 完成菲涅耳传播仿真，比较单焦点、双焦点和扩展焦深设计的轴向光强分布。 \\
第 4 周 & 根据最小线宽和制备条件修正设计参数，制作第一代二值振幅掩膜（先方案一，后方案二）。 \\
第 5 周 & 搭建激光准直和 CMOS 轴向扫描光路，完成方案一标准波带片聚焦与成像测试。 \\
第 6 周 & 系统测量不同形状及编码波带片的轴向光强、焦斑尺寸、旁瓣和有效焦深，并进行色差测试。 \\
第 7 周 & 进行成像离焦容忍度实验，比较不同设计在离焦条件下的图像质量。 \\
第 8 周 & 完成不确定度分析、实验与仿真对比、开题/中期汇报材料和最终展示图表。 \\
\bottomrule
\end{tabular}
\end{table}

% ===========================================================================
\section{希望向指导老师请教的问题}

为进一步提高方案可行性和科学性，拟向指导老师重点请教以下问题：

\begin{enumerate}[label=(\arabic*)]
\item 在给定最小线宽约为 $20\,\mu\mathrm{m}$ 的条件下，当前设计的孔径、焦距和波带数是否合理？
\item 高分辨率菲林、光绘掩膜、光刻或激光直写中，哪种方式最适合制作二值振幅型波带片？
\item 方案二-A 扇区编码结构中的角向边界会带来怎样的旁瓣和焦斑畸变？是否应增加扇区数或采用交错排列？
\item 方案二-B 径向二值优化中，目标函数应如何平衡目标匹配、旁瓣抑制和可制备性？
\item 如何更准确地测量最外圈线宽、边缘粗糙度和图案缩放误差，并将其纳入不确定度分析？
\item 若项目后续扩展，是否有可能将固定二值振幅波带片发展为相位型、可调型或光流控型波带片？
\end{enumerate}

% ===========================================================================
\section{结论}

本项目以菲涅耳波带片为研究对象，但不局限于传统的单焦点聚焦验证，而是将研究内容组织为两个递进的方案。方案一为制备约束下的稳妥基准方案，完整实现标准圆形波带片的设计、仿真、制备、聚焦、成像、色差与不确定度分析，并比较不同形状波带片的聚焦特性，确保项目满足竞赛基本要求并提供可靠对照。方案二为目标光场编码创新方案，将波带片重新理解为二值光场编码器，由目标轴向光强反向设计二值结构；其内部的方案二-A（确定性扇区编码）利用不同角向扇区对应不同焦距，实现直观的双焦点和扩展焦深，方案二-B（径向二值优化）将径向环带作为二值变量，通过目标函数优化主动逼近给定轴向光强分布。

两方案均可在普通物理光学实验平台上完成，不依赖复杂设备。其创新性主要体现在实验问题的重新定义、波带片设计范式由“固定焦距正向设计”向“目标光场驱动编码设计”的转变，以及对峰值强度、焦深、旁瓣和成像鲁棒性之间物理权衡的系统研究。该方案能够满足大学生物理实验竞赛对“设计实验方案、阐明工作原理、制作波带片、研究会聚能力、讨论适用范围和不确定度”的要求，同时具有较强的展示性和可拓展性。

\begin{thebibliography}{9}

\bibitem{goodman}
J. W. Goodman, \textit{Introduction to Fourier Optics}. Roberts and Company Publishers.

\bibitem{hecht}
E. Hecht, \textit{Optics}. Pearson.

\bibitem{bornwolf}
M. Born and E. Wolf, \textit{Principles of Optics}. Cambridge University Press.

\bibitem{kirz}
J. Kirz, ``Phase zone plates for x rays and the extreme uv,'' \textit{Journal of the Optical Society of America}, 1974.

\bibitem{doe}
相关衍射光学元件、双焦点波带片、扩展焦深波带片和光场优化设计文献可在后续正式报告中补充。

\end{thebibliography}

\end{document}
